\section{Finite descent and a complete elliptic-curve example}
\label{sec:descent}

Height-bounded enumeration becomes much stronger when a geometric transformation
strictly decreases an integer measure.  We now give an example in which descent
settles the entire rational-point problem, with no height cutoff left over.  The
main conclusions are Fermat's square-area theorem and the complete list of
rational points on
\[
  E_1:\qquad Y^2Z=X^3-XZ^2.
\]
The argument is elementary, but its logical implementation requires care:
``choose a least counterexample'' is replaced by a specified number of
primitive-recursive steps.  Fermat's result and its connection with the
congruent-number curves are classical; see \cite{ConradCongruent}.  An existing
Coq treatment of Fermat's theorem is useful background for formalization
\cite{DelahayeMayero}.  Neither reference is used as an unproved arithmetic
premise below.

\subsection{A descent principle with a numerical stopping bound}

\begin{lemma}[Finite descent in \iPRA]\label{lem:finite-descent}
Fix primitive-recursive functions $D$ and $\rho$ on codes of states, and a
primitive-recursive Boolean predicate $G$.  Suppose that \iPRA\ proves
\[
 G(s)\ \Longrightarrow\
  \bigl(0<\rho(s)\ \wedge\ G(D(s))\ \wedge\ \rho(D(s))<\rho(s)\bigr).
 \tag{D}
\]
Then \iPRA\ proves $\neg G(s)$ for every state $s$.
\end{lemma}

\begin{proof}
Define $s_0=s$ and $s_{j+1}=D(s_j)$ by primitive recursion.  Under the
assumption $G(s)$, quantifier-free induction on $j$ gives
\[
 G(s_j)\quad\hbox{and}\quad \rho(s_j)+j\leq\rho(s).
 \tag{1}
\]
Indeed, (D) gives $G(s_{j+1})$ and
$\rho(s_{j+1})+1\leq\rho(s_j)$, which proves the induction step by addition.
At $j=\rho(s)$, (1) implies $\rho(s_j)=0$, while (D) and $G(s_j)$ imply
$0<\rho(s_j)$.  This is a contradiction obtained from the assumption $G(s)$,
and therefore proves its negation intuitionistically.  There is no use of
double-negation elimination.  The displayed induction predicate is
quantifier-free after replacing the Boolean tests and inequalities by their
primitive-recursive characteristic functions.
\end{proof}

This is a schema for fixed primitive-recursive functions with their defining
equations.  It does not posit an internal universal evaluator for arbitrary
primitive-recursive programs.  In the application, the number of descent steps
is the original hypotenuse.  An abstract assertion that there is some smaller
counterexample would not by itself supply the required function $D$; the
formulas below do.

\subsection{Arithmetic tools without prime factorization}

All integers in the next three lemmas are nonnegative unless a sign is shown.
Division is used only after divisibility has been established.  Each exact
quotient can equivalently be implemented by a bounded search.

\begin{lemma}[The gcd tools]\label{lem:descent-gcd}
The following operations and implications have primitive-recursive witnesses
and quantifier-free-induction proofs.
\begin{enumerate}
\item If $a,b>0$, the greatest common divisor $g$ is a positive common divisor,
and there are signed integers $u,v$ with $ua+vb=g$.
\item If $\gcd(a,b)=1$ and $a\mid bc$, then $a\mid c$.
\item If $d^2\mid c^2$ and $d>0$, then $d\mid c$.
\item If $a,b,c>0$, $\gcd(a,b)=1$, and $ab=c^2$, then one can compute
positive integers $r,s$ such that $a=r^2$, $b=s^2$, and $c=rs$.
\end{enumerate}
\end{lemma}

\begin{proof}
For the first assertion, order the initial pair as
$(\max(a,b),\min(a,b))$ and perform Euclidean division.  Replace a pair
$(u,v)$ with $v>0$ by $(v,u\bmod v)$; after reaching a zero second entry,
leave the pair unchanged.  Before stopping, the sum of the entries strictly
decreases.  Thus $a+b+1$ iterations suffice.  Track the two entries as signed
linear combinations of the original $a,b$: subtraction of the quotient times
the second combination from the first gives the next combination.  The common
divisors are unchanged at each step, and the last nonzero entry divides both
preceding entries, hence both original entries.  The tracked last combination
is the required B\'ezout identity.  These statements are invariants of a fixed
finite computation and use quantifier-free induction.  If an explicit bound
for the coefficient storage is desired, each quotient is at most
$M=\max(a,b)$ and each coefficient is bounded in absolute value by
$(M+1)^{a+b+2}$.

For the second assertion, multiply a B\'ezout identity $ua+vb=1$ by $c$.
Both terms on the resulting right side are divisible by $a$ whenever
$a\mid bc$.  The same identity also shows that coprimality is preserved when
one of the entries is squared: squaring $ua+vb=1$ gives a B\'ezout identity
between $a$ and $b^2$.

For the third assertion, $c=0$ gives $d\mid c$ immediately. Otherwise
put $h=\gcd(c,d)$, $c=hr$, and $d=hs$.
Then $\gcd(r,s)=1$, and cancellation in $d^2\mid c^2$ gives $s^2\mid r^2$.
In particular $s\mid r^2$.  Applying the preceding coprimality observation
and the second assertion yields $s\mid1$, so $s=1$.  Consequently $d=h$ divides
$c$.

For the fourth assertion, put $h=\gcd(a,c)$, $a=hu$, and $c=hv$.
Here $\gcd(u,v)=1$, and the equation becomes
\[
 ub=hv^2.
\]
Since $u$ is coprime to $v^2$, it divides $h$.  Write $h=ue$.  Cancellation
now gives $a=e u^2$ and $b=e v^2$.  Thus $e$ divides both $a$ and $b$;
coprimality forces $e=1$.  Taking $r=u$ and $s=v$ proves the claim, including
$c=rs$.  This proof extracts the square roots by gcd and exact division.
It does not need a prime factorization of the inputs.
\end{proof}

Repeatedly applying the last assertion splits a square product of any fixed
number of pairwise coprime positive factors into squares.  To justify this
iteration, if $a$ is coprime to each of $b$ and $c$, multiply their B\'ezout
identities to obtain a B\'ezout identity for $a$ and $bc$.  Thus one factor is
coprime to the product of all the remaining factors.

\begin{lemma}[Normalization of square-area triangles]\label{lem:triangle-normalize}
From positive integers $a,b,c,s$ satisfying
\[
 a^2+b^2=c^2,\qquad ab=2s^2,
 \tag{2}
\]
one can compute positive integers $A,B,C,S$ satisfying the same equations
and $\gcd(A,B)=1$, with $C\leq c$.
\end{lemma}

\begin{proof}
Let $g=\gcd(a,b)$ and put $A=a/g$, $B=b/g$.  Since
$c^2=g^2(A^2+B^2)$, Lemma~\ref{lem:descent-gcd}(3) gives $g\mid c$;
put $C=c/g$.  The resulting Pythagorean equation is primitive in the two
legs.  Its legs cannot both be even, by coprimality, and cannot both be odd,
because their squared sum would be $2$ modulo $4$.  Thus exactly one leg is
even, and $K=AB/2$ is an integer.  Equation (2) now says $s^2=g^2K$.
Applying the same square-divisibility lemma gives $g\mid s$.  Put $S=s/g$.
Cancelling $g^2$ in the two equations proves the result.  All four outputs
are positive, and $C=c/g\leq c$.
\end{proof}

\begin{lemma}[Primitive Pythagorean parametrization]\label{lem:pythagorean-descent}
Let $A,B,C>0$, $A^2+B^2=C^2$, and $\gcd(A,B)=1$.  After exchanging the legs
so that $A$ is even, there are unique positive integers $m>n$ such that
\[
 A=2mn,\qquad B=m^2-n^2,\qquad C=m^2+n^2,
 \tag{3}
\]
where $\gcd(m,n)=1$ and $m,n$ have opposite parity.  They are obtained by
bounded arithmetic from $A,B,C$.
\end{lemma}

\begin{proof}
The parity argument in the preceding lemma shows that $B$ and $C$ are odd.
Also $\gcd(B,C)=1$: any common divisor $d$ satisfies $d^2\mid A^2$, hence
$d\mid A$ by Lemma~\ref{lem:descent-gcd}(3), and then $d=1$.
As $C>B>0$, the integers
\[
 U=\frac{C+B}{2},\qquad V=\frac{C-B}{2}
\]
are positive, with $U>V$.  A common divisor divides their sum $C$ and their
difference $B$, so $\gcd(U,V)=1$.  Their product is
\[
 UV=\frac{C^2-B^2}{4}=\left(\frac A2\right)^2.
\]
The square-splitting lemma gives $U=m^2$, $V=n^2$, and $A/2=mn$.
Addition and subtraction give (3), while $U>V$ gives $m>n$.
A common divisor of $m,n$ would divide $U,V$, so they are coprime.
Their squares sum to the odd number $C$, so they have opposite parity.
Finally, $m^2=(C+B)/2$ and $n^2=(C-B)/2$ determine the positive roots
uniquely: $x^2=y^2$ for nonnegative integers implies $x=y$ by
$(x-y)(x+y)=0$, or directly by strict monotonicity of squaring.
Each root is at most $C$ and can be found by bounded search.
\end{proof}

\subsection{The explicit descent}

\begin{theorem}[Fermat's square-area theorem in \iPRA]\label{thm:fermat-area}
There are no positive integers $a,b,c,s$ satisfying (2).  Consequently there
is no rational right triangle of positive side lengths whose area is a
nonzero rational square.  In particular, a rational right triangle cannot
have area $1$.
\end{theorem}

\begin{proof}
Normalize a supposed solution by Lemma~\ref{lem:triangle-normalize}.  We
will transform every primitive solution into another primitive solution
with a smaller hypotenuse.

Apply Lemma~\ref{lem:pythagorean-descent} and substitute (3) into the
square-area equation.  Cancelling $2$ gives
\[
 m n(m+n)(m-n)=S^2.                         \tag{4}
\]
These four factors are positive and pairwise coprime.  For pairs involving
$m$ or $n$, a common divisor would also divide both $m$ and $n$.
For the remaining pair, a common divisor of $m+n$ and $m-n$ divides
$2m$ and $2n$, and hence divides $2$ by a B\'ezout identity for $m,n$.
Both $m+n$ and $m-n$ are odd, so that common divisor is $1$.
The square-splitting lemma applied three times to (4) supplies positive
integers $r,s,t,u$ with
\[
 m=r^2,\qquad n=s^2,\qquad m+n=t^2,\qquad m-n=u^2.
 \tag{5}
\]
In particular $t,u$ are odd, $t>u>0$, and $\gcd(t,u)=1$.
Moreover $s$ is even.  Indeed $m,n$ have opposite parity; if $m$ were even
and $n$ odd, then the first two equations of (5) would give
$m-n\equiv-1\equiv3\pmod4$, contradicting $m-n=u^2$.
Thus $n$ is even, and $n=s^2$ forces $s$ even.

Define
\[
 A'=\frac{t+u}{2},\qquad B'=\frac{t-u}{2},\qquad
 C'=r,\qquad S'=\frac s2.                  \tag{6}
\]
All four are positive integers.  Direct calculation using (5) gives
\begin{align*}
 (A')^2+(B')^2
   &=\frac{t^2+u^2}{2}=m=r^2=(C')^2,\\
 A'B'&=\frac{t^2-u^2}{4}=\frac n2
       =\frac{s^2}{2}=2(S')^2.
\end{align*}
A common divisor of $A',B'$ divides their sum $t$ and difference $u$,
so the new triangle is primitive.  Finally,
\[
 0<C'=r\leq r^2=m<m^2+n^2=C.
 \tag{7}
\]
This completes the explicit descent step.

For its \iPRA\ implementation, let $G$ be the Boolean predicate that a
quadruple is positive, satisfies (2), and has coprime legs.  On a valid input,
let $D$ carry out the leg exchange, the bounded square-root extractions in
(3) and (5), and formula (6).  On an invalid input let $D$ return the zero
quadruple.  Bounded integer square root, gcd, exact quotient, comparison,
and tuple coding are primitive recursive, so this defines a total
primitive-recursive function.  The preceding paragraphs prove (D) with
$\rho(a,b,c,s)=c$.  Lemma~\ref{lem:finite-descent} therefore proves that
$G$ is always false.  Normalization then rules out (2) without the
coprimality condition.

For the rational statement, suppose positive rationals $a,b,c$ satisfy
$a^2+b^2=c^2$ and $ab/2=q^2$ with $q\ne0$.  Take the product $L>0$ of the
positive denominators of $a,b,c,q$.  The numbers
$La,Lb,Lc,L|q|$ are positive integers satisfying (2), a contradiction.
The product of denominators is a specified primitive-recursive operation;
no least common denominator is needed.
\end{proof}

The proof has a useful complexity audit.  On a valid state of hypotenuse
$C$, the parameter searches need inspect no integer exceeding $C$.
The next state's coordinates are all at most $C$: the two legs are smaller
than the new hypotenuse, and its square-area root is smaller than that
hypotenuse as well.  Iteration for exactly the original $C$ steps is
therefore more than sufficient.  The contradiction concerns the purported
input itself, so it does not require an a priori bound on all possible
counterexamples.

\begin{corollary}[A quartic obstruction]\label{cor:quartic-obstruction}
There are no positive integers $x>y$ and $z$ with $x^4-y^4=z^2$.
In particular $a^4+b^4=c^4$ has no solution in positive integers.
\end{corollary}

\begin{proof}
A solution of the first equation would give the integer right triangle
\[
 \bigl(2x^2y^2,\ x^4-y^4,\ x^4+y^4\bigr).
\]
The identity $(2uv)^2+(u^2-v^2)^2=(u^2+v^2)^2$, with $u=x^2$ and
$v=y^2$, verifies that it is right-angled.  Its area is
\[
 x^2y^2(x^4-y^4)=(xyz)^2>0,
\]
contrary to Theorem~\ref{thm:fermat-area}.  For the second assertion,
$c^4-a^4=(b^2)^2$ would contradict the first, since $c>a>0$.
\end{proof}

\begin{corollary}[Three rational squares]\label{cor:square-progression}
There are no rational numbers $r,s,t,d$, with $d\ne0$, for which
$s^2-r^2=t^2-s^2=d^2$.
\end{corollary}

\begin{proof}
Put $R=|r|$, $S=|s|$, $T=|t|$.  Then $T>R\geq0$ and $S>0$.
The positive rational numbers
\[
  T-R,\quad T+R,\quad 2S
\]
form a right triangle because
$(T-R)^2+(T+R)^2=2(T^2+R^2)=4S^2$.
Its area is $(T^2-R^2)/2=d^2>0$, again a contradiction.
\end{proof}

\subsection{A rational correspondence with cubic curves}

For a positive rational $n$, write
\[
 E_n:\quad y^2=x^3-n^2x.
\]
Its projective equation is $Y^2Z=X^3-n^2XZ^2$; after clearing the denominator
of $n$, this is an integral homogeneous equation in the format used earlier.
We first keep signs in the triangle coordinates.  This makes the formulas
mutually inverse without dividing into geometric sign cases.

\begin{theorem}[Congruent-number correspondence]\label{thm:En-correspondence}
For each positive rational $n$, there are mutually inverse rational maps
between the following two collections of rational tuples:
\[
 \{(a,b,c):a^2+b^2=c^2,\ ab=2n\}
 \quad\longleftrightarrow\quad
 \{(x,y):y^2=x^3-n^2x,\ y\ne0\}.
\]
They are
\begin{align}
 (a,b,c)&\longmapsto
   \left(\frac{nb}{c-a},\frac{2n^2}{c-a}\right),
   \label{eq:triangle-to-curve}\\
 (x,y)&\longmapsto
   \left(\frac{x^2-n^2}{y},\frac{2nx}{y},
                     \frac{x^2+n^2}{y}\right).
   \label{eq:curve-to-triangle}
\end{align}
Every indicated denominator is nonzero.  These are primitive-recursive maps
on reduced rational codes, with any fixed default output outside their stated
domains.  Thus $E_n$ has a rational point with $y\ne0$ if and only if there
is a positive rational right triangle of area $n$.
\end{theorem}

\begin{proof}
If $(a,b,c)$ is in the first collection, then $ab=2n\ne0$, so $a,b$ are
nonzero.  Also $c-a\ne0$: otherwise $b^2=c^2-a^2=0$ and hence $b=0$.
Set $\delta=c-a$, and define $x,y$ by
\eqref{eq:triangle-to-curve}.  Then $y\ne0$.  The two identities
\[
 b^2-\delta^2=2a\delta,\qquad
 b^2+\delta^2=2c\delta
\]
follow by expanding and using $a^2+b^2=c^2$.  They imply
\[
 x^2-n^2=ay,\qquad x^2+n^2=cy,
 \qquad 2nx=by.
 \tag{8}
\]
In addition $ax=anb/\delta=2n^2/\delta=y$.  Multiplying the first
identity in (8) by $x$ gives
$x(x^2-n^2)=axy=y^2$, so the image lies on $E_n$.
The three identities in (8) show that applying
\eqref{eq:curve-to-triangle} recovers exactly $a,b,c$.

Conversely, let $(x,y)$ lie on $E_n$ with $y\ne0$.  Since
$y^2=x(x^2-n^2)$, both $x$ and $x^2-n^2$ are nonzero.  Define $a,b,c$
by \eqref{eq:curve-to-triangle}.  The identity
\[
 (x^2-n^2)^2+(2nx)^2=(x^2+n^2)^2
\]
gives $a^2+b^2=c^2$, and
\[
 ab=\frac{2nx(x^2-n^2)}{y^2}=2n.
\]
Furthermore $c-a=2n^2/y\ne0$, so substitution into
\eqref{eq:triangle-to-curve} returns $x,y$.  This proves both the domain
checks and the inverse assertions.

For positivity, $ab=2n>0$ makes $a,b$ nonzero with the same sign; also
$c^2=a^2+b^2>0$.  Replacing $a,b,c$ by their absolute values therefore
produces positive lengths and preserves both equations.  Conversely a
positive triangle is already a tuple in the first collection.  All sign
tests and rational operations are decidable and primitive recursive on
integer numerator and positive denominator codes.
\end{proof}

As an exact numerical check, the triangle $(3,4,5)$ has area $6$ and gives
$(x,y)=(12,36)$ on $E_6$, since $36^2=12^3-36\cdot12$.  Formula
\eqref{eq:curve-to-triangle} sends it back to $(3,4,5)$.

\begin{theorem}[Complete rational points of $E_1$]\label{thm:E1}
In the coding of rational projective points, \iPRA\ proves that
\[
 E_1(\Q)=\{[0:1:0],\ [0:0:1],\ [1:0:1],\ [-1:0:1]\}.
 \tag{9}
\]
The braces in (9) abbreviate the finite disjunction that every rational
projective point satisfying the cubic equation is equal to one of the four
listed points, together with the direct verification that all four are on
the cubic.
Projective notation permits either sign; the last displayed point has
canonical integer code $(1,0,-1)$ under our first-nonzero-positive convention.
\end{theorem}

\begin{proof}
Let $[X:Y:Z]$ satisfy $Y^2Z=X^3-XZ^2$.
The equality $Z=0$ is decidable.  In that case $X^3=0$, so $X=0$, and
nonzeroness of the coordinate tuple gives $Y\ne0$.  The point is therefore
$[0:1:0]$.

If $Z\ne0$, put $x=X/Z$ and $y=Y/Z$.  They satisfy
$y^2=x^3-x$.  The equality $y=0$ is again decidable.  If $y\ne0$,
Theorem~\ref{thm:En-correspondence} produces a positive rational right
triangle of area $1$, contradicting Theorem~\ref{thm:fermat-area}.
Thus $y=0$, and $x(x-1)(x+1)=0$.  Rational equality is decidable and
$\Q$ has no zero divisors: test $x=0$, and if it is false cancel $x$;
then test $x-1=0$, and if it is false cancel $x-1$.  This yields
$x=0$, $x=1$, or $x=-1$.  They give exactly the three affine points in
(9).  Substitution verifies each of the four displayed projective points.
\end{proof}

For the usual geometric interpretation, the cubic is nonsingular.  With
$F=Y^2Z-X^3+n^2XZ^2$, its partial derivatives are
\[
 F_X=-3X^2+n^2Z^2,\qquad F_Y=2YZ,\qquad
 F_Z=Y^2+2n^2XZ.
\]
At the unique point at infinity, $F_Z\ne0$.  In the affine chart, simultaneous
vanishing of $F$ and $F_Y$ gives $y=0$ and $x\in\{0,n,-n\}$.
At these points $F_X$ is respectively $n^2,-2n^2,-2n^2$, all nonzero.
The calculation is valid over every characteristic-zero field with $n\ne0$.
It identifies $E_n$ as a nonsingular plane cubic with a rational distinguished
point, the customary elliptic-curve situation.  The point classification
itself is entirely the integer and rational calculation above; it requires
no group-law or general finite-generation theorem.

\subsection{What the formalization must establish}

The proof naturally separates into five reusable Lean components:
\begin{enumerate}
\item The signed-integer and rational-code algebra, including bounded gcd,
exact division, and the four implications in
Lemma~\ref{lem:descent-gcd}.
\item Primitive Pythagorean normalization and the parameter extraction
of Lemmas~\ref{lem:triangle-normalize} and
\ref{lem:pythagorean-descent}.
\item A total function \texttt{descendSquareArea} on quadruple codes, with
its preservation and strict-decrease theorems proved under the Boolean
validity hypothesis.
\item A general finite-iteration lemma using only the displayed
quantifier-free invariant, followed by
\texttt{noSquareAreaTriangle}.
\item The two rational transformations, denominator proofs, inverse
identities, and the finite disjunction \texttt{rationalPointsE1}.
\end{enumerate}

The included executable reference file \texttt{code/descent.py} implements
the arithmetic transformations and checks their identities on exact rational
examples.  It also implements the total descent map.  A valid input for its
main descent branch cannot exist, by the theorem; therefore finite testing
cannot exercise that branch on an actual valid counterexample or replace
its mathematical proof.  The branch's preservation proof is the symbolic
calculation (4)--(7).  The remaining formalization task is to check those
symbolic calculations and their arithmetic preliminaries in Lean, and then
to connect the fixed algorithms and the indicated quantifier-free
inductions to the chosen object-theory representation of \iPRA.
